\section{Introduction and Data Sources}

We built a PostgreSQL database linking HDB resale transactions to nearby hawker centres~\cite{datagovHDB,datagovHawker}.
The database supports queries on resale prices, flat attributes, remaining lease and hawker proximity.
This report explains our schema, data processing, six SQL queries and performance experiments on PostgreSQL~18.

\subsection{Objectives}
The project has five objectives.
\begin{enumerate}
  \item design a proximity-aware ER model and a constrained seven-table schema;
  \item acquire and transform official open data into one CSV per table;
  \item answer required queries Q1--Q5 plus an original Q6 that joins every table;
  \item measure client-side runtime as the database grows from one to five resale files;
  \item select and evaluate indexes on the largest database.
\end{enumerate}

\subsection{Datasets}
Table~\ref{tab:datasets} lists the inputs used by the revised ETL.
Resale files span January~1990 through September~2026; earlier files record approval date and later files registration date~\cite{datagovHDB}.
Hawker attributes come from the NEA government list~\cite{datagovHawker}; coordinates come from course-provided \texttt{hawker\_geo.csv}~\cite{hawkerGeo}.
HDB block coordinates come primarily from \texttt{hdb.csv}~\cite{siteselectHDB}, with OneMap as a cached fallback~\cite{onemap}.

\begin{table}[H]
\centering
\caption{Source datasets used in this project.}
\label{tab:datasets}
\begin{tabularx}{\textwidth}{@{}p{3.1cm}p{4.2cm}X@{}}
\toprule
Source & Role & Coverage / notes \\
\midrule
HDB resale (5 CSV files) & Transactions, flat attributes, towns & 1990--2026; 987{,}686 sales at full load \\
NEA hawker list & Centre type, stall counts, address & 107 government centres \\
\texttt{hawker\_geo.csv} & Hawker coordinates & From NTULearn GeoJSON extract \\
\texttt{hdb.csv} & Building coordinates & BlueSkyLT/\texttt{siteselect\_sg} \\
OneMap Search API & Fallback geocoding & Cached under \texttt{data/cache/} \\
\bottomrule
\end{tabularx}
\end{table}

\subsection{Pipeline overview}
We use the same ETL code to build five cumulative datasets, then load each into the seven-table schema.
Q1--Q6 are measured at each size before indexes are evaluated on the largest database (Figure~\ref{fig:pipeline}).
Executable artefacts are listed in Appendix~\ref{app:artefacts}.

\begin{figure}[H]
\centering
\fbox{\parbox{0.92\textwidth}{\centering\small
Raw CSVs $\rightarrow$ Python ETL $\rightarrow$ 7 table CSVs $\rightarrow$
\texttt{schema.sql} + \texttt{COPY} $\rightarrow$ PostgreSQL\\[0.4em]
$\rightarrow$ Q1--Q6 + timing (5 sizes) $\rightarrow$ index screening $\rightarrow$ final \texttt{CREATE INDEX}}}
\caption{End-to-end pipeline implemented in the Hexa Nexus repository.}
\label{fig:pipeline}
\end{figure}
